\section{Introduction}
\label{sec:introduction-v128}

A rank condition records where multiplication fails.  The scheme
structure of that rank condition records how it fails.  In the
cube-zero multiplication problem considered here this distinction is
geometric: the reduction is a Schubert divisor, whereas successive
nilpotent layers recover a polarized Jacobian K3 surface, contact
divisors on a Segre quadric, vertex-primary corrections, and the
higher-rank boundary on which those objects specialize.

The central point of this paper is that these layers are not artifacts
of a chosen determinantal presentation.  They are powers of the
intrinsic nilradical of the failure scheme.  The determinantal
presentation is used to calculate them; the resulting graded schemes,
their associated points, and their multiplication are intrinsic.

Revision 128 has three purposes.  First, it proves a determinant
completion theorem that applies to every symmetric power and contains
the global depth-six theorem of the \((1,4,6)\) problem as the case
\((e,r)=(4,2)\).  Second, it makes the corank-two atlas genuinely
generic: the Hilbert and primary tables are computed over rational
function fields on transverse stabilizer slices, rather than inferred
from isolated integral examples.  Third, it proves the relative
associated-point statement needed to pass from finite coefficient
matrices to fibrewise primary signatures and supplies explicit
one-parameter boundary degenerations.

\subsection{The multiplication failure scheme}

Let \(V\) be a complex vector space of dimension
\(e\in\{3,4\}\), put \(p=e(e-1)/2\), and let
\[
 R\subset\Sym^2V
\]
be an \(e\)-dimensional relation space.  Set
\[
 S_R=\Sym^2V/R,\qquad
 B_R=\C\oplus V\oplus S_R,
\]
with multiplication induced by the quotient
\(\Sym^2V\twoheadrightarrow S_R\) and with
\((V\oplus S_R)S_R=0\).  On
\[
 X_R=\Gr(e,V\oplus S_R)
\]
let \(\mathcal K\) be the tautological subbundle.  For every
\(m\ge2\) the multiplication map defines the same zeroth-Fitting
scheme
\begin{equation}\label{eq:failure-intro}
 \widehat D_R=
 V\!\left(
 \Fitt_0\coker[
 \Sym^m(\OO\oplus\mathcal K)\longrightarrow
 B_R\otimes\OO]\right).
\end{equation}
Its reduction is the Schubert divisor
\(\widehat\Delta_R\) on which \(\mathcal K\to V\) drops rank.

Regard \(R\) as a web of quadrics on \(\Pj(V^*)\).  On the
invariant open \(G_e^\circ\) used throughout, the web is
basepoint-free and its Jacobian hypersurface
\begin{equation}\label{eq:Y-intro}
 Y_R=V\!\left(\det(\partial q_i/\partial x_j)\right)
 \subset\Pj(V^*)
\end{equation}
is smooth.  For \(e=4\) it is a quartic K3 surface.

\subsection{Intrinsic reconstruction}

On the projection-corank-at-most-one locus let \(D_R\) be the
restriction of \(\widehat D_R\), let \(\NN_R\) be its
nilradical, and let \(E_R\) be the annihilator scheme of the first
nilpotent layer.  The first main theorem is the reconstruction result
proved in Section~\ref{sec:polarization}.

\begin{theorem}[Polarized reconstruction]
\label{thm:main-reconstruction}
For \(R\in G_e^\circ\), put
\[
 \LL_R=\omega_{E_R}\otimes
 \NN_R^{\otimes-(p+e-1)}.
\]
Then
\[
 \bigoplus_{n\ge0}H^0(E_R,\LL_R^{\otimes n})
 \simeq
 \bigoplus_{n\ge0}H^0(Y_R,\OO_{Y_R}(en))
\]
as graded algebras.  In particular, for \(e=4\) the abstract
failure scheme determines the polarized quartic K3
\((Y_R,\OO_{Y_R}(1))\).
\end{theorem}

\subsection{The intrinsic graded boundary packet}

Write \(\widehat\NN_R\) for the nilradical of
\(\OO_{\widehat D_R}\).  The schemes
\[
 \ZZ_j(\widehat D_R)=
 V\!\left(\Ann(\widehat\NN_R^{,j})\right)
\]
and the products among
\(\widehat\NN_R^j/\widehat\NN_R^{j+1}\) are invariant under
every abstract isomorphism of \(\widehat D_R\).  On a Schur chart
the reduced equation is \(d=\det T\) and
\[
 \mathcal I_{\widehat D_R}=dJ.
\]
The residual-colon theorem of Section~\ref{sec:stratified-nilpotent}
identifies
\begin{equation}\label{eq:intro-graded-packet}
 \widehat\NN_R^j/\widehat\NN_R^{j+1}
 \simeq
 \mathcal L^{\otimes j}\otimes\OO_{W_j},\qquad
 W_j=V\bigl((d)+(J:d^{j-1})\bigr).
\end{equation}
The coordinate expression changes from chart to chart; the modules
and schemes in \eqref{eq:intro-graded-packet} do not.

\begin{theorem}[Boundary structure]
\label{thm:main-depth}
For \(R\in G_4^\circ\) the following hold.
\begin{enumerate}
\item On every projection-rank chart,
\[
 d^5\in J,\qquad \widehat\NN_R^6=0.
\]
This is the \((e,r)=(4,2)\) case of the universal determinant
completion theorem, whose exponent
\(\binom{e+r-1}{r-1}\) is sharp in the universal class.

\item The multiplication of the positive graded pieces is the
canonical quotient multiplication
\[
 \OO_{W_i}\otimes\OO_{W_j}\longrightarrow\OO_{W_{i+j}}
\]
with the expected conormal-line twist.  No additional scalar
multiplication constants occur.

\item If the pure block \(A_H\) is injective at projection corank
two, the mixed block has nine orbit types.  On the generic
coefficient stratum of every orbit, the complete nonzero graded
supports through \(W_3\), their Hilbert functions, radicals,
associated supports, and vertex-primary corrections are given in
Theorems~\ref{thm:corank-two-mixed-kernel-atlas} and
\ref{thm:generic-primary-signatures}.  Genericity is proved over
rational function fields on transverse stabilizer slices.

\item If \(A_H\) drops rank at projection corank two, every
generic rank type \((a,b)\) has an explicit primary signature.
Besides the Segre quadric, the only positive-dimensional generic
supports that occur are reduced divisors of two or four lines in one
ruling.  The remaining layers are vertex-primary; in the
\((a,b)=(1,3)\) row a length-three embedded vertex accompanies
the two ruling lines.

\item At projection corank three, \(d^3\notin J\) and
\(d^5\in J\).  The two possible depths are separated by the
intrinsic membership condition \(d^4\in J\).  The rank-drop
passage in the representation argument is functorial, not inferred
from semicontinuity.

\item At projection corank four,
\[
 d^3\notin J,qquad d^4\in J.
\]
The second membership is proved from the actual Cauchy--Pieri
multiplication in \(\C[\End(V)]\): the product of the
\((4,4,4,0)\)- and \((4)\)-components has nonzero
\((4,4,4,4)\)-projection.  Since that line is spanned by
\(d^4\), the local nilpotency index is exactly five.
\end{enumerate}
\end{theorem}

Thus the phrase \emph{boundary atlas} refers to intrinsic graded
data, not to the chosen coefficient matrices.  The Macaulay matrices
form a finite presentation atlas.  Any two such presentation atlases
have a common refinement and compute the same intrinsic primary
signature.

\subsection{Primary specialization}

The preceding finite calculations are promoted to family statements
in Section~\ref{sec:relative-primary-specialization}.  A crucial
order of operations is enforced there.  First one stratifies so that
the finite presentations defining the colons and their cokernels are
flat; only then is kernel formation pulled to fibres.  Relative
assassins are subsequently spread out and the base is refined by
Noetherian induction.  This gives constant fibrewise associated
points, embedded/minimal status, incidences, and generic lengths.

The same section writes explicit one-parameter degenerations of the
mixed block.  For example an off-Segre kernel point specializes to a
Segre point; a secant kernel line specializes to a tangent or either
ruling line; rank-two kernels specialize to the two rank-one types;
and both rank-one types specialize to the zero mixed map.  The
graded schemes in the two fibres are computed, so the specialization
records the birth of vertex-primary and ruling components rather than
merely changing orbit labels.

\subsection{The simple-contact locus}

The general atlas contains the earlier primary theorem as its split
form.  We retain the following statement because it is the geometric
entry point to the collision analysis.

\begin{theorem}[Corank-two primary structure on the simple-contact open]
\label{thm:main-coranktwo}
There is a nonempty invariant open
\(G_4^\dagger\subset G_4^\circ\) such that, for each
\(R\in G_4^\dagger\), a dense open set of rank-two images
\(H\subset V\) has, after an etale extension and smooth auxiliary
coordinates,
\[
 \mathcal I_{\widehat D_R}
 =
 (\delta)\cap P_1^2\cap P_2^2\cap P_3^2\cap P_4^2
 \cap\mathfrak m^5.
\]
The rank-two stratum is an additional associated support and
\(\widehat\NN_R^3=0\ne\widehat\NN_R^2\).
\end{theorem}

Repeated roots, the containment locus, the decomposable-kernel wall,
and all remaining mixed-kernel orbits are treated without splitting
the contact quartic.

\subsection{Variation and finite ambiguity}

\begin{theorem}[Variation and finite ambiguity]
\label{thm:main-moduli}
For \(e=4\), the image of
\[
 R\longmapsto(Y_R,\OO_{Y_R}(1))
\]
in quartic K3 moduli has dimension nine.  On a nonempty invariant
open the induced map from projective equivalence classes of relation
spaces to its image is generically finite.
\end{theorem}

The nine-dimensional count is classical Reye geometry, not a novelty
claim.  The manuscript-specific inverse statement is the recovery of
the polarized K3 from an abstract nonreduced failure scheme.  Revision
128 does not need an unproved generic-injectivity assertion: its
broader structural consequence is instead the universal symmetric-power
determinant-completion theorem.

\subsection{Organization and novelty boundary}

Section~\ref{sec:relations} fixes the relation model and
Section~\ref{sec:fitting} gives the Fitting presentation.
Section~\ref{sec:depth} develops the intrinsic nilpotent filtration,
followed by polarized reconstruction.  Section
\ref{sec:boundary-atlas-statement} states the finite boundary theorems;
the inherited corank-two and global residual calculations follow, and
Section~\ref{sec:boundary-atlas-proofs} proves the new determinant,
genericity, and higher-corank statements.  Section
\ref{sec:relative-primary-specialization} proves the family and
specialization results.  Deformation to the normal cone is then used
only as the standard package for the computed graded algebra.  The
moduli and classical/priority comparisons close the article.

Deformation to the normal cone, flattening stratification, Schur
functors, and standard monomial straightening are classical tools.
The new assertions are the specific residual Fitting algebra, the
universal determinant-completion mechanism applied to it, effective
Pieri nonvanishing for the corank-four membership, the generic and
primary corank-two boundary tables, their specialization laws, and
the reconstruction of the polarized K3 from the intrinsic nilpotent
scheme.  Ballico's 1993 paper remains cited as a documentary boundary;
no theorem-level nonanticipation claim is made without its complete
text.
